% Abhakseite „Das kann ich" – Zone, Einheit 1, Einheit 2
\begin{abhakseite}
\abhakgruppe{Kennst du schon}
\abhak{1}{Ich kann den Prozentsatz berechnen.}
\abhak{2}{Ich kann Prozentangaben als Dezimalzahl schreiben.}
\abhak{3}{Ich kann die relative Häufigkeit angeben.}
\abhak{4}{Ich kann Werte aus einem Säulendiagramm ablesen.}
\abhak{5}{Ich kann einen Winkel als Anteil vom Vollkreis angeben.}
\abhak{6}{Ich kann einen Winkel mit dem Geodreieck zeichnen.}
\abhak{7}{Ich kann Längen in Zentimetern und Millimetern abtragen.}
\abhak{8}{Ich kann Kenngrößen einer Liste bestimmen.}
\abhak{9}{Ich kann einen Bruchteil und ein Vielfaches einer Zahl berechnen.}
\abhak{10}{Ich finde den Fehler beim Anteil.}
\abhak{11}{Ich kann einen Anteil an allen berechnen.}
\abhakgruppe{Einheit 1 · Streifen- und Kreisdiagramm}
\abhak{12}{Ich erkenne, was das Ganze ist.}
\abhak{13}{Ich erkenne, ob eine Zahl ein Anteil oder ein Winkel ist.}
\abhakgruppe{Streifendiagramm}
\abhak{14}{Ich kann einen Anteil als Streifenabschnitt zeichnen.}
\abhak{15}{Ich kann ein Streifendiagramm mit mehreren Abschnitten zeichnen.}
\abhakgruppe{Kreisdiagramm}
\abhak{16}{Ich kann den Mittelpunktswinkel zu einem Anteil berechnen.}
\abhak{17}{Ich kann einen Sektor ab dem Radius zeichnen.}
\abhak{18}{Ich kann ein Kreisdiagramm zeichnen. (kennst du wahrscheinlich seit Klasse 6)}
\abhak{19}{Ich kann den Anteil eines Sektors schätzen.}
\abhak{20}{Ich kann Sektoren nach ihrer Größe zuordnen.}
\abhak{21}{Ich kann Sektoren nach ihrer Größe zuordnen – weiter.}
\abhak{22}{Ich finde den Fehler beim Kreisdiagramm.}
\abhak{23}{Ich kann begründen, warum 1\,\% im Kreisdiagramm 3,6° sind.}
\abhak{24}{Ich kann Umfrageergebnisse in einem Kreisdiagramm darstellen.}
\abhakgruppe{Einheit 2 · Diagramme beurteilen und Boxplot}
\abhak{25}{Ich erkenne, wo die Achse anfängt.}
\abhakgruppe{Aussagen zu einem Diagramm prüfen}
\abhak{26}{Ich kann eine Aussage zu einem Diagramm prüfen.}
\abhak{27}{Ich kann Aussagen mit Anteilen zu einem Diagramm prüfen.}
\abhakgruppe{Verzerrte Diagramme erkennen}
\abhak{28}{Ich kann erklären, warum eine abgeschnittene Achse täuscht.}
\abhak{29}{Ich kann erklären, warum eine gedehnte Achse täuscht.}
\abhak{30}{Ich kann die Fortschreibung eines Trends prüfen. (kommt in Klasse 9)}
\abhak{31}{Ich kann die Fortschreibung eines Trends prüfen – weiter.}
\abhakgruppe{Boxplot}
\abhak{32}{Ich kann einen Boxplot lesen. (kommt in Klasse 8, je nach Buch bis Kl. 9)}
\abhak{33}{Ich kann einen Boxplot zeichnen.}
\abhak{34}{Ich kann zwei Boxplots vergleichen.}
\abhak{35}{Ich finde den Fehler bei der Deutung eines Diagramms.}
\abhak{36}{Ich kann begründen, warum die Achse eines Säulendiagramms bei 0 beginnen sollte.}
\abhak{37}{Ich kann Diagramme in der Werbung beurteilen.}
\end{abhakseite}
